\section{Introduction}
\label{sec:introduction}

In a UTXO ledger, a payment spends an output that an earlier transaction
created. If the recipient must wait until that transaction executes
publicly, every hop of a payment chain adds a confirmation delay.
Consider Alice paying Bob 100 CAL, with Bob forwarding the received
output to Carol. To forward this output before public execution, Bob needs verifiable
evidence that its value is backed. If Bob's payment then executes first, the ledger has
consumed an output that Alice's transaction has not yet created, and the
missing 100 CAL needs an identified source of funds. The accounting must
also stay correct whether Alice's transaction arrives before Bob's, after
it, or never.

Certificate-based payment systems already separate recipient-side
usability from public processing. FastPay and Zef support low-latency
certified payments, and Sui Lutris combines broadcast and consensus paths
for different kinds of state access~\cite{fastpay,zef,lutris}.
Payment-channel protocols reach rapid transfers through channel
relationships and their liquidity~\cite{blitz,donner}. Next addresses how certified payments are funded on a native ledger when the source of a spent output has not yet executed. A certified output whose source
has not yet executed should support a publicly executed successor, and an
identified party with funds set aside should answer for the missing
source. Users spend their own routed outputs, registered organizations
back unresolved sources, and counterparties need no prefunded channel
path.

Four coupled problems arise. First, certificates circulate before public
registration. Accounting that starts from publicly recorded gaps misses commitments that
already exist, so signing capacity must cover the outstanding CAL liability
of every certified payment, including payments never published. Second, a source may arrive after the
committee has compensated its missing value. The late arrival must not pay
the same value both to the guarantor that advanced it and to the recipient
who already spent it, or delay itself would move funds. Third, a client
may hold the material of a source and never submit it. The issuer's own members should be able to resubmit an available source
before compensation becomes due, while compensation should proceed
without waiting for the historical adaptation service. Fourth, once a source has been compensated, a reader
of the stored history should be able to see which reserve funded the
consumed input. Recording that fact must preserve the owner
authorization, the output identifiers, and the commitments that
successors reference. A stable transaction hash alone does not specify
or authorize such a change.

Next fixes payment rights and assigns unresolved sources to their direct
issuers. An output certificate (TXCer) binds a payment's outputs to an
identified issuing organization. After verifying and atomically recording
an output and its TXCer, the recipient requests a successor through the
output's designated consumption organization. The request carries certificates for its direct inputs. Owner authorization, spend
state, fees, and resource limits are checked separately.

If a certified input is still missing from public creation state, the
committee atomically consumes it, records a direct obligation against the
issuer, and creates the child outputs. The child outputs are final while
the obligation is open. The block that carries the child also carries the
certificate of the payment that created the consumed output. Members of
the issuing organization that approved that payment rebuild it from their
stored approvals and resubmit it, without new signatures. If the source
executes, the obligation is fulfilled. If it remains absent when a compensation decision executes after its
deadline, that decision debits the issuer's reserve,
which supplies the missing backing, and the recipient receives no second
credit. The decision needs no threshold adaptation. A source that arrives
later repays the compensating issuer in the same transition and creates no
second output for the recipient. When all sources eventually execute,
delay changes only the temporary exposure of reserves, and the final
allocation of principal is the same as without delay.

The compensation decision is itself a public record that closes the
obligation. Representation commands then replace the funding reference of the
consumed input with the reserve reference, under the existing commitments.
The stored block names the reserve debit that funded the input and still
verifies against its original block identity, so a reader can check the
funding source of a designated historical input in the block itself.
Wallets, synchronization, and replay continue to read original commands.
We organize the design into three contributions with distinct
responsibilities.

\noindent\textbf{C1: A layered architecture for one-round continuation.}
We separate member certification, recipient delivery, evidence
replication, and public execution. A member locks inputs and reserves
capacity atomically before it signs. One parallel certification round produces a TXCer. INSTALL replicates
complete material within the organization, while public submission
proceeds independently; both run in the background. Registered organizations use one evidence interface within and
across organizations, and payment principal stays separate from guarantor
reserves.

\noindent\textbf{C2: Direct responsibility with funded, delay-neutral
accounting.} We define the atomic missing-input transition and its closure
paths: fulfillment by the source, compensation by a public decision, and
repayment by a late source. For fixed input and certificate shapes, the
ancestry evidence of a request does not grow with transfer depth. Under
the fixed-membership Byzantine model, we prove unique consumption and
asset conservation, and we prove that signing capacity bounds the
outstanding CAL liability of certificates, including unpublished ones. For a
finite set of non-conflicting payments whose sources all execute, any
interleaving of arrivals and compensations ends with the same CAL holdings
and reserve balances as undelayed execution.

\noindent\textbf{C3: Evidence-driven source resolution with
identity-preserving representation.} We let the original signers resubmit a source from the certificate that
its successor publishes, and separate compensation decisions from threshold
adaptation. Restricted representation commands have no economic effect and
preserve owner authorizations, output references, and the complete block
identity, including its part-set header. Commands targeting one block share
adaptation of the affected parts' final representations, while installation
advances directly to a committed revision. For a fixed set of decided slots,
the final body does not depend on how the commands are grouped.

We evaluate two builds on one M4 Max host: the continuation and accounting
baseline, and the decision-first build that adds resubmission and the
separate compensation decision. On the baseline, 100-hop chains reach
their last certified receipt in a three-run median of 158.846\,ms, against
65.204\,s with public-confirmation waits between hops. Three
30,000-payment runs complete 1,750--1,887 payments/s including background
drain, and a nine-run matrix closes 21,600 payments and repays all 216
compensations. On the decision-first build, three-run functional tests show
a withheld source executing without compensation while the recipient's
receipt takes about 14\,ms. Compensation and repayment complete with all
adaptation services paused, and the per-run receipt P50 and P95 ranges overlap in a six-run
ordinary-path comparison.
Sections~\ref{sec:model} and~\ref{sec:protocol} define the design, and
Sections~\ref{sec:security} and~\ref{sec:evaluation} develop its safety
argument and evaluation.
